\documentclass[11pt,a4paper]{article}

% ---- packages -----------------------------------------------------------------
\usepackage[margin=2.4cm]{geometry}
\usepackage{amsmath,amssymb}
\usepackage{booktabs}
\usepackage{array}
\usepackage{tabularx}
\usepackage{xcolor}
\usepackage{graphicx}
\usepackage{enumitem}
\usepackage{microtype}

% colours (define BEFORE hyperref so its options resolve)
\definecolor{NavyBlue}{RGB}{30,80,150}
\definecolor{RuleGrey}{RGB}{120,120,120}
\definecolor{BoxBack}{RGB}{244,247,251}
\definecolor{BoxLine}{RGB}{200,212,228}

\usepackage[hidelinks,colorlinks=true,linkcolor=NavyBlue,urlcolor=NavyBlue,citecolor=NavyBlue]{hyperref}

% section styling
\usepackage{sectsty}
\sectionfont{\color{NavyBlue}\large}
\subsectionfont{\color{NavyBlue}\normalsize}

% headers/footers
\usepackage{fancyhdr}
\pagestyle{fancy}
\fancyhf{}
\renewcommand{\headrulewidth}{0.3pt}
\fancyhead[L]{\footnotesize\color{RuleGrey} Conflict $\to$ Road-Blockage Framework}
\fancyhead[R]{\footnotesize\color{RuleGrey} ACLED Afghanistan 2017--2026}
\fancyfoot[C]{\thepage}

% a light "insight/note" box
\usepackage[most]{tcolorbox}
\newtcolorbox{notebox}{colback=BoxBack,colframe=BoxLine,boxrule=0.6pt,
  arc=2pt,left=6pt,right=6pt,top=5pt,bottom=5pt}

% code font for inline identifiers
\newcommand{\code}[1]{\texttt{\small #1}}

% math shortcuts
\newcommand{\Reff}{R_{\mathrm{eff}}}
\newcommand{\Rphys}{R_{\mathrm{phys}}}
\newcommand{\Rgeo}{R_{\mathrm{geo}}}
\newcommand{\Msev}{M_{\mathrm{sev}}}
\newcommand{\Pblock}{P_{\mathrm{block}}}

\title{\vspace{-1.2cm}\textbf{\color{NavyBlue} A Framework for Estimating Road-Blockage
Probability from Conflict Events}\\[2pt]
\large A supervised, parametric impact model}
\author{Conflict-event $\to$ Road-Blockage (\textsc{c2rb}) model}
\date{ACLED, Afghanistan, 2017--2026 \quad(69{,}655 geocoded events)}

\begin{document}
\maketitle
\thispagestyle{fancy}
\vspace{-0.4cm}

% =====================================================================================
\section{Problem and approach}

\textbf{Goal.} Given a conflict event (or a set of events over a time window), estimate which
nearby roads are \emph{affected} and the \emph{probability that each road is blocked}. The
central modelling question is: \textbf{how large is an event's impact ``buffer zone'', and how
does it differ by \code{sub\_event\_type}?}

\textbf{Ground-truth labels.} Each event carries two human/LLM-verified yes/no labels:
\code{is\_road\_affected} (the event impacted a road in any way --- 20.8\% of events) and
\code{is\_road\_blocked} (a road was actually closed / made impassable --- 0.87\% of events); the
latter is our calibration and scoring target. \emph{(The framework was originally built
weakly-supervised --- the signal was mined from the \code{notes} free text --- because the raw
export had no such column; that proxy has now been replaced by these labels.)}

\textbf{Our approach --- a supervised parametric probability model.} We build an explicit,
interpretable probability model whose parameters are (a)~\textbf{seeded from published
literature} (radii and decay shapes) and (b)~\textbf{calibrated against the ground-truth
\code{is\_road\_blocked} label}. This is defensible, reproducible, and mirrors established
practice: ACLED's own \emph{Conflict Exposure} product uses fixed event-type buffers
(1/2/5\,km) with Voronoi de-duplication, and a published conflict-impact model assigns
per-event-type radii (Battles 25\,km, Explosions 10\,km, Protests/Riots 5\,km) with
type-specific decay shapes.

\begin{notebox}
\textbf{Pipeline.}\quad
$\text{event}\;\longrightarrow\;\Reff\;\longrightarrow\;K(d)\;\longrightarrow\;
\Pblock(d)\;\xrightarrow{\;\text{noisy-OR}\;+\;\text{time decay}\;}\;P(\text{road blocked})$
\end{notebox}

% =====================================================================================
\section{The model}

For an event $e$ with sub-event-type $s$, geo-precision code $g\in\{1,2,3\}$, fatalities
$n_f$, and civilian-targeting flag $c\in\{0,1\}$:

\subsection{Effective buffer radius --- ``how big is the zone?''}

\begin{equation}
\Reff(e) \;=\; \sqrt{\,\Rphys(s)^2 + \Rgeo(g)^2\,}\;\cdot\;\Msev(e)
\label{eq:reff}
\end{equation}

\begin{itemize}[nosep,leftmargin=1.4em]
  \item $\Rphys(s)$ --- physical/operational impact radius of the event type (km),
        literature-seeded per \code{sub\_event\_type}. E.g.\ IED 1\,km, armed clash 15\,km,
        protest 4\,km.
  \item $\Rgeo(g)$ --- \emph{locational uncertainty} of where the event actually happened.
        ACLED geo-precision $1\approx$ town (1\,km), $2\approx$ near a town / part of a region
        (5\,km), $3\approx$ wider region (25\,km). A low-precision event must spread its
        blockage probability over a larger uncertain area.
  \item \textbf{Quadrature combination} $\sqrt{\Rphys^2+\Rgeo^2}$. Both terms behave like
        independent Gaussian length-scales (physical spread and locational error); variances of
        independent Gaussians add, so their standard deviations add in quadrature. This is more
        principled than naive addition (which double-counts) and never smaller than either term.
\end{itemize}

The severity multiplier $\Msev(e)$ stretches the radius for deadlier and civilian-directed
events, with diminishing returns so a 50-fatality battle does not produce a $50\times$ radius:
\begin{align}
S(n_f)   &= \alpha\,(1+n_f) + (1-\alpha)\bigl(1+\ln(1+n_f)\bigr) \label{eq:sev}\\[2pt]
\Msev(e) &= \operatorname{clip}\!\Bigl(\Bigl(1 + \kappa\Bigl(\tfrac{S(n_f)}{S(n_{\mathrm{ref}})}-1\Bigr)\Bigr)
            \,\cdot\,(1+\gamma_{\mathrm{civ}}\,c),\;0.5,\,3.0\Bigr) \label{eq:msev}
\end{align}
with defaults $\alpha=0.15$ (mostly log growth), $\kappa=0.25$, $\gamma_{\mathrm{civ}}=0.30$.
The \code{clip} bounds keep outliers from exploding the buffer.

\paragraph{Explanation of Equation~\eqref{eq:sev} (the severity score $S(n_f)$).}
A deadlier event tends to affect a larger area, but not linearly --- the jump from $0$ to $5$
fatalities matters far more than $100$ to $105$. Equation~\eqref{eq:sev} captures this by
\textbf{blending two growth curves}:
\begin{itemize}[nosep,leftmargin=1.4em]
  \item the \textbf{linear term} $(1+n_f)$ grows in direct proportion to the fatality count ---
        it lets genuinely large, mass-casualty events widen the buffer; and
  \item the \textbf{logarithmic term} $\bigl(1+\ln(1+n_f)\bigr)$ grows ever more slowly as $n_f$
        rises --- it applies \emph{diminishing returns} so extreme counts do not blow up the
        radius without bound.
\end{itemize}
The mixing weight $\alpha\in[0,1]$ sets the balance: $\alpha=0$ is pure-log (gentle), $\alpha=1$
is pure-linear (aggressive); the default $0.15$ keeps growth mostly logarithmic with a small
linear component. The $1+\,$offsets are deliberate --- they make $\ln(1+n_f)$ well-defined at
$n_f=0$ (avoiding $\ln 0=-\infty$) and guarantee $S(0)=1$, so a zero-fatality event is the
natural baseline. $S$ enters Equation~\eqref{eq:msev} only through the \emph{ratio}
$S(n_f)/S(n_{\mathrm{ref}})$: the event is compared to a reference fatality level
$n_{\mathrm{ref}}$ (default $1$), the ratio is turned into a gentle multiplier by $\kappa$,
bumped by $\gamma_{\mathrm{civ}}$ if civilians were targeted, and the product is bounded to
$[0.5,3]$. In
short, Equation~\eqref{eq:sev} converts a raw body count into a \emph{bounded,
diminishing-returns severity signal} that stretches the buffer sensibly rather than mechanically.

\subsection{Distance-decay blockage probability --- ``how does it fall with distance?''}

\begin{equation}
\Pblock(d,e) \;=\; P_0(s)\,\cdot\,K_s\!\left(d;\,\Reff\right)
\label{eq:pblock}
\end{equation}

\begin{itemize}[nosep,leftmargin=1.4em]
  \item $P_0(s)$ --- \emph{peak blockage propensity}: the probability that an event of type $s$
        blocks a road at the epicentre. \textbf{This is the parameter calibrated on the
        \code{is\_road\_blocked} label} (Sec.~\ref{sec:calib}).
  \item $K_s(d)\in[0,1]$ --- distance-decay kernel, its \emph{shape} chosen by the event
        type's physics (Table~\ref{tab:kernels}).
\end{itemize}

\begin{table}[h]
\centering
\renewcommand{\arraystretch}{1.35}
\begin{tabularx}{\textwidth}{@{}l l l X@{}}
\toprule
\textbf{Kernel} & \textbf{Formula} & \textbf{Used for} & \textbf{Rationale} \\
\midrule
Gaussian & $\exp\!\left(-\dfrac{d^2}{2\sigma^2}\right),\ \sigma=\Reff/2$
  & IED, suicide, grenade, shelling, airstrike & sharp, blast-like, fast decay \\
Exponential & $\exp\!\left(-\dfrac{d}{\lambda}\right),\ \lambda=\Reff/3$
  & armed clash / battles & heavier tail --- fighting spreads along terrain/roads \\
Uniform & $\mathbf{1}\{d\le \Reff\}$
  & protests, riots, territorial control & roughly constant over an area, then cut off \\
Point & $\mathbf{1}\{d\le \max(\Reff,0.3)\}$
  & targeted violence, arrests & effect confined to the site \\
\bottomrule
\end{tabularx}
\caption{Distance-decay kernels, chosen per event type.}
\label{tab:kernels}
\end{table}

\subsection{Combining many events --- noisy-OR with temporal decay}

A road segment near several events is blocked if \emph{any} of them blocks it. Treating events
as independent blocking attempts gives the \textbf{noisy-OR}, with an \textbf{exponential
temporal decay} $\varphi(\Delta t)=0.5^{\,\Delta t/\tau}$ (half-life $\tau=30$ days) so older
events fade:
\begin{equation}
P_{\text{segment}} \;=\; 1 - \prod_{e\,\in\,W}\Bigl[\,1 - \Pblock(d_e,e)\,\cdot\,\varphi(\Delta t_e)\,\Bigr]
\label{eq:noisyor}
\end{equation}
over a time window $W$. ACLED-style Voronoi de-duplication can be layered on to avoid
double-counting tightly-clustered events (noted as an extension).

\subsection{From points to roads}

We use the WFP/AGCHO Afghanistan road network
(\code{afg\_trs\_roads\_l\_arazi\_2023x\_fixed\_segmented.shp}): \textbf{27{,}991 LineString
segments, $\approx$1\,km each} (median 1{,}068\,m), EPSG:4326, each carrying road name, class,
average slope and altitude. Each segment is represented by its midpoint; $d_e$ is the
segment-to-event distance. Every segment within $r_{\max}=3\,\Reff$ of an event receives that
event's $\Pblock(d_e,e)$, and segments are aggregated by noisy-OR. The result is a per-segment
blockage probability that rolls up to \textbf{named roads}. Without a road layer the same
scoring runs over a regular lat/lon grid, producing a continuous blockage-risk raster (the
fallback) --- the code path is identical, only the target points differ.

% =====================================================================================
\section{Calibration \& validation on the ground-truth labels}
\label{sec:calib}

\subsection{The two labels --- and why we model \emph{blocked}, not \emph{affected}}

Each event carries two yes/no labels: \code{is\_road\_affected} (14{,}456 events, 20.8\%) and
\code{is\_road\_blocked} (607 events, 0.87\%). Blockage is almost a strict subset of
affectedness --- only \textbf{1 of 607} blocked events is not also flagged affected --- so the
labels are internally consistent ($P_0^{\text{blk}}\le P_0^{\text{aff}}$ by construction) and
$P(\text{blocked}\mid\text{affected})=0.042$. The two tell \emph{different stories}, which is
exactly why we calibrate on \emph{blocked}:

\begin{table}[h]
\centering
\renewcommand{\arraystretch}{1.2}
\begin{tabular}{@{}l r r r@{}}
\toprule
\textbf{sub\_event\_type} & \textbf{$n$} & \textbf{P(affected)} & \textbf{P(blocked)} \\
\midrule
Remote explosive/landmine/IED & 8{,}545 & \textbf{0.656} & \textbf{0.003} \\
Suicide bomb                  & 289     & 0.405 & 0.000 \\
Violent demonstration         & 44      & 0.386 & \textbf{0.182} \\
Peaceful protest              & 1{,}454 & 0.102 & 0.061 \\
Air/drone strike              & 6{,}249 & 0.049 & 0.001 \\
\bottomrule
\end{tabular}
\caption{Affected $\neq$ blocked: an IED \emph{affects} roads constantly (roadside bomb) but
rarely \emph{blocks} them; a demonstration is the reverse.}
\label{tab:divergence}
\end{table}

An IED/landmine \emph{affects} roads constantly (66\% of the time) but rarely \emph{blocks} them
(traffic resumes --- 0.3\%); a protest or demonstration deliberately \emph{blocks} the road. A
single ``is the road relevant?'' signal conflates these; the ground-truth \code{is\_road\_blocked}
label separates them cleanly.

\subsection{Peak blockage propensity $P_0(s)$}

$P_0(s)=P(\code{is\_road\_blocked}=\text{yes}\mid s)$ is estimated with \textbf{Beta--Binomial
empirical-Bayes shrinkage} (prior strength 50, anchored at the 0.87\% global rate) so
small-sample types (Grenade $n{=}176$, Suicide $n{=}289$) get stable estimates instead of
noisy zeros:
\begin{equation}
\hat{P}_0(s) \;=\; \frac{k_s + a_0}{n_s + a_0 + b_0},
\qquad a_0 = \bar p\,m,\quad b_0=(1-\bar p)\,m,\quad m=50,
\end{equation}
where $k_s$ positives out of $n_s$ events of type $s$, and $\bar p$ is the global rate.
Selected results with Beta posterior 95\% credible intervals (full 24-row table in
\code{countries/afghanistan.yaml}):

\begin{table}[h]
\centering
\renewcommand{\arraystretch}{1.2}
\begin{tabular}{@{}l r r l@{}}
\toprule
\textbf{sub\_event\_type} & \textbf{$n$} & \textbf{$P_0$ (shrunk)} & \textbf{95\% CrI} \\
\midrule
Violent demonstration                 & 44       & \textbf{0.090} & [0.04, 0.16] \\
Change to group/activity              & 318      & 0.083 & [0.06, 0.11] \\
Non-violent transfer of territory     & 130      & 0.080 & [0.05, 0.13] \\
Peaceful protest                      & 1{,}454  & 0.060 & [0.05, 0.07] \\
Non-state actor overtakes territory   & 905      & 0.060 & [0.05, 0.08] \\
Armed clash                           & 39{,}974 & 0.007 & [0.006, 0.008] \\
Remote explosive/landmine/IED         & 8{,}545  & 0.003 & [0.002, 0.005] \\
Air/drone strike                      & 6{,}249  & \textbf{0.001} & [0.000, 0.002] \\
\bottomrule
\end{tabular}
\caption{Calibrated blockage propensity --- the headline result.}
\label{tab:p0}
\end{table}

\begin{notebox}
\textbf{Headline result.} Road \emph{blockage} is driven by \textbf{protests and
territorial-control events} (deliberate, sustained closures), not by explosive violence --- the
\emph{opposite} of an affected-based ranking. This confirms that the model \textbf{must} vary by
\code{sub\_event\_type}, and shows why the ground-truth label matters: it corrects the old weak
label, which wrongly placed IEDs at the top by conflating \emph{affected} with \emph{blocked}.
\end{notebox}

\subsection{Validation --- supervised, on the true label}

A logistic regression $\code{is\_road\_blocked}\sim\code{sub\_event\_type}+\log(1+\text{fatalities})
+\code{geo\_precision}+\code{civilian\_targeting}$, evaluated with 5-fold cross-validation:

\begin{itemize}[nosep,leftmargin=1.4em]
  \item \textbf{AUC $= 0.73$} --- event attributes discriminate road-blocking events well above
        chance.
  \item \textbf{PR-AUC $= 0.042$} vs a 0.0087 prevalence --- $\sim$5$\times$ better than random
        under heavy imbalance.
  \item \textbf{Brier $= 0.0085$} --- \emph{better} than the no-skill baseline (0.0086), i.e.\
        genuinely informative.
  \item \textbf{Well-calibrated} --- predicted $\approx$ observed across deciles. Use the plain
        model for calibrated probabilities; \code{class\_weight="balanced"} improves ranking but
        inflates probabilities (apply isotonic/Platt calibration if class-weighting).
\end{itemize}

The model is both discriminative and calibrated, so the $P_0(s)$ values are trustworthy as
probabilities, not merely rankings.

\subsection{Other sanity checks (in the notebook)}
\begin{itemize}[nosep,leftmargin=1.4em]
  \item \textbf{Roads-within-1\,km benchmark --- passed.} With the real road network loaded,
        \textbf{78\% of events fall within 1\,km of a road}, independently matching the
        published $\sim$70\% figure --- strong external validation.
  \item \textbf{Nesting check} --- blocked $\subseteq$ affected (1 exception in 607), confirming
        label consistency.
  \item \textbf{Sensitivity analysis} over $\Rphys$ and decay-kernel choice shows how the
        affected-road footprint responds to assumptions (what is assumed vs.\ learned).
\end{itemize}

% =====================================================================================
\section{Statistical methods used}

Distance-decay kernels (Gaussian / exponential / uniform); quadrature combination of
independent uncertainty scales; bounded log-linear severity scaling; noisy-OR probabilistic
aggregation with exponential temporal decay; Beta--Binomial empirical-Bayes shrinkage with Beta
posterior credible intervals (and Wilson intervals on raw rates); supervised logistic regression
with cross-validated AUC / PR-AUC / Brier / calibration-curve diagnostics under class imbalance;
(optional) Voronoi tessellation for overlap de-duplication and KDE for hotspot context.

% =====================================================================================
\section{Limitations \& honest caveats}

\begin{enumerate}[leftmargin=1.6em]
  \item \textbf{Labels are ground truth, but derived from \code{notes}.} \code{is\_road\_affected}
        / \code{is\_road\_blocked} were classified from the ACLED \code{notes} text, so they
        inherit its reporting quality --- an event whose notes omit a closure can be a false
        negative. \textbf{Class imbalance is severe} (blocked $=0.87\%$), so absolute
        probabilities are small and PR-style metrics matter more than accuracy. Cross-checking
        against independent closure data (OCHA/Logistics Cluster access reports) would further
        harden $P_0$ --- the calibration code is unchanged.
  \item \textbf{$\Rgeo$ cannot be learned from this data.} ACLED snaps all events at a named
        \code{location} to one shared centroid, so intra-location coordinate spread is
        structurally $\sim$0\,km. We therefore fall back to ACLED's \emph{documented} precision
        semantics (1/5/25\,km). With finer coordinates, the empirical-spread estimator in the
        notebook would populate $\Rgeo$.
  \item \textbf{$\Rphys$ and kernel shapes are assumptions} (literature-seeded), unlike $P_0$
        (data-calibrated). The sensitivity analysis quantifies how much they matter; tune in
        \code{params.base.yaml}.
  \item \textbf{Independence in noisy-OR} slightly over-counts spatially correlated events;
        Voronoi de-duplication (ACLED's approach) is the recommended refinement.
  \item \textbf{Static / impact-only.} This scores the impact of given events; it does not
        forecast where future events will occur.
\end{enumerate}

% =====================================================================================
\section{Pipeline, configuration, and localization}

The framework runs headless through three CLI scripts, all reading a \emph{layered} configuration.
\code{python train.py --config countries/<name>.yaml} calibrates $P_0(s)$ on the country's
\code{is\_road\_blocked} labels (Beta--Binomial shrinkage), writes it back into the country config,
and emits \code{artifacts/<name>/} (\code{model.joblib}, \code{metrics.json}, \code{model\_card.md},
calibration plots). Training produces \emph{two models for two questions}: the calibrated
parametric $P_0$ (the interpretable core that drives the spatial scorer), and a logistic-regression
classifier that is cross-validated for the \S3.3 metrics and then refit on all rows and persisted
into \code{model.joblib} as the deployable event-level model.
\code{python score.py --config countries/<name>.yaml --as-of <date> --window-days <n>}
scores road segments via noisy-OR over the window's events and ranks the most likely-blocked named
roads, falling back to a grid surface if no road layer is set --- the \emph{spatial} question
(\emph{which roads} are blocked). For the \emph{event-level} question (\emph{will this event} block
a road?), \code{python predict.py --config countries/<name>.yaml --sub-event-type \ldots} loads the
deployable \code{model.joblib} --- the classifier \code{train.py} fits on the labels and persists
alongside $P_0$ --- and returns one probability per event (single event or CSV batch).

Country-independent physics --- $\Rphys$, decay kernels, $\Rgeo$, the severity and temporal
constants --- lives in \code{params.base.yaml} and is shared by every country. Each country adds a
small override (\code{countries/<name>.yaml}) carrying only its data paths and calibrated $P_0$
(plus an optional \code{crs\_metric} hint for GIS preprocessing); \code{c2rb.load\_config} merges
the two. Adapting the model to a new country means copying the override, pointing it at that
country's ACLED export and segmented road network, and re-running \code{train.py} then
\code{score.py} --- no changes to the \code{c2rb} package. See \code{docs/COUNTRY\_GUIDE.md} for
the procedure (including \code{tools/segment\_roads.py}, which cuts a raw road network into
$\sim$1\,km segments).

\begin{table}[h]
\centering
\renewcommand{\arraystretch}{1.2}
\begin{tabularx}{\textwidth}{@{}l X@{}}
\toprule
\textbf{File} & \textbf{Purpose} \\
\midrule
\code{c2rb/} & Framework core package: geometry, calibration, config, model physics, spatial scorer, classifier \\
\code{train.py / score.py / predict.py} & CLI: calibrate + save the model / per-road rankings (spatial) / event-level P(blocked) \\
\code{app.py} & Streamlit app: event-level prediction + click-to-place spatial scoring \\
\code{artifacts/<name>/model.joblib} & Deployable trained model: event-level classifier + parametric $P_0$ \\
\code{params.base.yaml} & Universal parameters: $\Rphys$, decay class, $\Rgeo$, severity \& time constants (shared) \\
\code{countries/<name>.yaml} & Per-country config: data paths, calibrated $P_0$ (written by \code{train.py}) \\
\code{docs/} & Stage-by-stage pipeline documentation (data $\to$ labeling $\to$ training $\to$ evaluation $\to$ deployment) \\
\code{docs/COUNTRY\_GUIDE.md} & Step-by-step guide to building a model for another country \\
\code{tests/ / Makefile} & Pytest smoke suite over the core math and config; make targets for the pipeline \\
\code{tools/segment\_roads.py} & Utility: cut a raw road network into $\sim$1\,km segments \\
\code{conflict\_road\_blockage.ipynb} & Annotated end-to-end walkthrough: load $\to$ label $\to$ calibrate $\to$ score $\to$ map $\to$ sensitivity \\
\code{build\_notebook.py} & Regenerates the notebook from cell definitions (diffable source of truth) \\
\code{requirements.txt} & Pinned dependencies for reproducibility \\
\code{methodology.tex / .pdf} & This document \\
\bottomrule
\end{tabularx}
\end{table}

\textbf{Reproduce.} Create the venv and \code{pip install -r requirements.txt}, then either run the
CLI scripts above or execute the notebook top-to-bottom. The road and admin layers are already
wired in \code{countries/afghanistan.yaml}; set \code{paths.roads: null} to fall back to
grid-raster mode.

% =====================================================================================
\section{References}

\begin{itemize}[nosep,leftmargin=1.4em]
  \item ACLED, \emph{Conflict Exposure methodology} (buffer zones 1/2/5\,km, event-type
        adjustment, Voronoi de-duplication).\\
        \url{https://acleddata.com/methodology/conflict-exposure-methodology}
  \item ACLED, \emph{Codebook} (event types, sub-event types, geo-precision codes) ---
        \code{AcledCodebook.md}.
  \item \emph{Applying Computational Engineering Modelling to Analyse the Social Impact of
        Conflict and Violent Events} (per-type impact radii \& decay shapes), PMC12562815.\\
        \url{https://pmc.ncbi.nlm.nih.gov/articles/PMC12562815/}
  \item \emph{Proximity to roads shapes patterns of violence in North and West Africa}
        ($\sim$70\% of violent events within 1\,km of a road).\\
        \url{https://anl.geog.ufl.edu/roads-violence-nwafrica/}
  \item Road network: WFP / AGCHO-NSIA Afghanistan segmented roads
        (\code{afg\_trs\_roads\_l\_arazi\_2023x}); admin boundaries: AGCHO-NSIA Afghanistan
        adm1/adm2 (2017).
\end{itemize}

\end{document}
